% Chapter 10: Gamma Function
% ch10_gamma.tex

\chapter{The Gamma Function and Euler's Reflection Formula}

\begin{historical}
Euler introduced the Gamma function in 1729 as an interpolating function for the factorial,
communicating it to Goldbach in letters (Opera Omnia, Series I, Vol. 14). He discovered
the integral representation, the infinite product, and the reflection formula, establishing
the Gamma function as a cornerstone of analysis. The Gamma function generalizes the factorial
to complex numbers (except non-positive integers) and appears throughout mathematics, physics,
and statistics. \cite{euler1729}\end{historical}

\section{Definition}

\begin{definition}[Euler's Integral of the Second Kind]
For $\Re(z) > 0$:
\[
\Gamma(z) = \int_0^\infty t^{z-1} e^{-t}\, dt
\]
\end{definition}

\begin{theorem}[Factorial Interpolation]
For any positive integer $n$:
\[
\Gamma(n+1) = n!
\]
\end{theorem}

\begin{proof}
Integrating by parts with $u = t^z$ and $dv = e^{-t}dt$:
\[
\Gamma(z+1) = \left[ -t^z e^{-t} \right]_0^\infty + z \int_0^\infty t^{z-1} e^{-t}\, dt = 0 + z\Gamma(z) = z\Gamma(z).
\]
Since $\Gamma(1) = \int_0^\infty e^{-t}dt = 1$, induction yields $\Gamma(n+1) = n!$.
\end{proof}

\begin{lstlisting}[caption={Gamma function values}, label=lst:gamma]
from euler import gamma_function, digamma_function
import math

print("Gamma function values \Gamma(n+1) = n!:")
for n in range(10):
    result = gamma_function(n + 1)
    print(f"  \Gamma({n+1}) = {result:.6f}  (n! = {math.factorial(n)})")

print("\nGamma at half-integers:")
for x in [0.5, 1.5, 2.5, 3.5]:
    result = gamma_function(x)
    print(f"  \Gamma({x}) = {result:.6f}")
\end{lstlisting}

\section{Fundamental Properties}

\subsection{Recurrence Relation}

\begin{theorem}[Functional Equation]
For all $z \in \C \setminus \{0, -1, -2, \ldots\}$:
\[
\Gamma(z+1) = z \Gamma(z).
\]
\end{theorem}

\begin{proof}
Integration by parts on the integral definition gives:
\[
\Gamma(z+1) = \int_0^\infty t^z e^{-t} dt = \left[-t^z e^{-t}\right]_0^\infty + z\int_0^\infty t^{z-1} e^{-t} dt = z\Gamma(z).
\]
\end{proof}

\subsection{Analytic Continuation}

\begin{theorem}[Meromorphic Extension]
The Gamma function extends to a meromorphic function on $\mathbb{C}$ with simple poles at $z = 0, -1, -2, \ldots$. The residue at $z = -n$ is $\frac{(-1)^n}{n!}$.
\end{theorem}

\begin{proof}
Using the recurrence $\Gamma(z) = \frac{\Gamma(z+1)}{z}$, we can extend $\Gamma$ leftward. At $z = -n$, the function has a simple pole with residue:
\[
\lim_{z \to -n} (z + n)\Gamma(z) = \frac{(-1)^n}{n!}.
\]
\end{proof}

\section{Weierstrass Product Representation}

Euler's original formulation used an infinite product:
\[
\frac{1}{\Gamma(z)} = z e^{\gamma z} \prod_{n=1}^\infty \left(1 + \frac{z}{n}\right) e^{-z/n}
\]
where $\gamma$ is Euler's constant. This product converges for all complex $z$.

\begin{theorem}[Weierstrass Product]
The Weierstrass form of the Gamma function is:
\[
\Gamma(z) = \frac{e^{-\gamma z}}{z} \prod_{n=1}^\infty \left(1 + \frac{z}{n}\right)^{-1} e^{z/n}.
\]
\end{theorem}

\begin{proof}
Taking the logarithm of the product and differentiating, one can verify that the resulting function satisfies the same functional equation as $\Gamma(z)$ and agrees with it at positive integers.
\end{proof}

\section{Key Formulas}

\subsection{Euler's Reflection Formula}

\begin{theorem}[Euler's Reflection Formula]
\[
\Gamma(z)\Gamma(1-z) = \frac{\pi}{\sin(\pi z)}
\]
\end{theorem}

\begin{proof}
Using the Weierstrass infinite product representation of the Gamma function:
\[
\frac{1}{\Gamma(z)} = z e^{\gamma z} \prod_{n=1}^\infty \left(1 + \frac{z}{n}\right) e^{-z/n},
\]
we have:
\[
\frac{1}{\Gamma(z)\Gamma(-z)} = -z^2 \prod_{n=1}^\infty \left(1 - \frac{z^2}{n^2}\right).
\]
Using the relation $\Gamma(1-z) = -z\Gamma(-z)$, this simplifies to:
\[
\frac{1}{\Gamma(z)\Gamma(1-z)} = z \prod_{n=1}^\infty \left(1 - \frac{z^2}{n^2}\right) = \frac{\sin(\pi z)}{\pi},
\]
which yields $\Gamma(z)\Gamma(1-z) = \frac{\pi}{\sin(\pi z)}$.
\end{proof}

\begin{corollary}
Setting $z = 1/2$ gives $\Gamma(1/2)^2 = \pi$, so $\Gamma(1/2) = \sqrt{\pi}$.
\end{corollary}

\subsection{Legendre Duplication Formula}

\begin{theorem}[Legendre Duplication Formula]
\[
\Gamma(z)\Gamma\left(z+\frac{1}{2}\right) = 2^{1-2z}\sqrt{\pi}\,\Gamma(2z)
\]
\end{theorem}

\begin{proof}
This follows from the Gauss multiplication formula with $m = 2$:
\[
\Gamma(z)\Gamma\left(z + \frac{1}{2}\right) = \frac{\sqrt{\pi}}{2^{2z-1}} \Gamma(2z).
\]
\end{proof}

\subsection{Gauss Multiplication Formula}

\begin{theorem}[Gauss Multiplication Formula]
For any positive integer $m$:
\[
\prod_{k=0}^{m-1} \Gamma\left(z + \frac{k}{m}\right) = \frac{(2\pi)^{(m-1)/2}}{m^{z-1/2}} \Gamma(mz).
\]
\end{theorem}

\section{The Beta Function}

\begin{definition}[Beta Function]
The Beta function is defined as:
\[
B(x, y) = \int_0^1 t^{x-1}(1-t)^{y-1} dt = \frac{\Gamma(x)\Gamma(y)}{\Gamma(x+y)}, \qquad \Re(x) > 0, \Re(y) > 0.
\]
\end{definition}

\begin{theorem}[Beta-Gamma Relation]
\[
B(x, y) = \frac{\Gamma(x)\Gamma(y)}{\Gamma(x+y)}.
\]
\end{theorem}

\begin{proof}
Starting from the integral representations:
\[
\Gamma(x)\Gamma(y) = \int_0^\infty \int_0^\infty s^{x-1} t^{y-1} e^{-(s+t)}\, ds\, dt.
\]
Substituting $s = uv$, $t = u(1-v)$ gives:
\[
\Gamma(x)\Gamma(y) = \int_0^1 v^{x-1}(1-v)^{y-1} dv \int_0^\infty u^{x+y-1} e^{-u} du = B(x,y) \Gamma(x+y).
\]
\end{proof}

\section{The Digamma Function}

The digamma function is the logarithmic derivative:
\[
\psi(z) = \frac{d}{dz} \log \Gamma(z) = \frac{\Gamma'(z)}{\Gamma(z)}
\]

\begin{lstlisting}[caption={Digamma function}, label=lst:digamma]
from euler import digamma_function
import math

print("Digamma function \psi(z):")
for z in [1, 2, 3, 4, 0.5, 1.5]:
    val = digamma_function(z)
    print(f"  \psi({z}) = {val:.8f}")
\end{lstlisting}

\begin{theorem}[Series Representation]
For $\Re(z) > 0$:
\[
\psi(z) = -\gamma + \sum_{n=0}^\infty \left(\frac{1}{n+1} - \frac{1}{n+z}\right).
\]
\end{theorem}

\begin{theorem}[Special Values]
\[
\psi(1) = -\gamma, \quad \psi(1/2) = -\gamma - 2\ln 2, \quad \psi(n) = -\gamma + H_{n-1}
\]
where $H_n$ is the $n$-th harmonic number.
\end{theorem}

\section{Stirling's Formula}

\begin{theorem}[Stirling's Formula]
As $z \to \infty$ in the right half-plane:
\[
\Gamma(z) \sim \sqrt{2\pi} \, z^{z-1/2} e^{-z} \left(1 + \frac{1}{12z} + \frac{1}{288z^2} - \cdots\right).
\]
\end{theorem}

\begin{proof}
Using the integral representation and Laplace's method, the dominant contribution comes from near the maximum of the integrand at $t = z$.
\end{proof}

\subsection{Refined Stirling Series}

A more accurate approximation is:
\[
\Gamma(z) \approx \sqrt{\frac{2\pi}{z}} \left(\frac{z}{e}\right)^z \left(1 + \frac{1}{12z} + \frac{1}{288z^2} - \frac{139}{51840z^3} - \cdots\right).
\]

\section{Computational Methods}

\subsection{Lanczos Approximation}
For numerical computation of $\Gamma(z)$:
\[
\Gamma(z+1) = \sqrt{2\pi} \left(z + g + \frac{1}{2}\right)^{z+1/2} e^{-(z+g+1/2)} \sum_{n=0}^{N} \frac{c_n}{z+n}
\]
where $g$ is a tuning parameter and $c_n$ are coefficients.

\subsection{Spouge's Approximation}
An alternative to Lanczos with better error control:
\[
\Gamma(z) \approx \sqrt{2\pi}(z+a)^{z+1/2}e^{-(z+a)}\left(c_0 + \sum_{n=1}^N \frac{c_n}{z+n}\right)
\]

\section{Applications}

\begin{itemize}
\item \textbf{Probability:} Beta function $B(x,y) = \frac{\Gamma(x)\Gamma(y)}{\Gamma(x+y)}$
\item \textbf{Statistics:} Gamma distribution, Chi-squared distribution
\item \textbf{Combinatorics:} Counting permutations with repetition
\item \textbf{Number theory:} Functional equation of the Riemann zeta function
\item \textbf{Physics:} Partition functions in statistical mechanics
\item \textbf{Quantum mechanics:} Wavefunction normalization
\end{itemize}

\begin{remark}[Gamma Distribution]
The gamma distribution has PDF:
\[
f(x; \alpha, \beta) = \frac{\beta^\alpha}{\Gamma(\alpha)} x^{\alpha-1} e^{-\beta x}, \qquad x > 0.
\]
\end{remark}

\begin{remark}[Chi-Squared Distribution]
If $X_1, \ldots, X_n$ are i.i.d. standard normal, then $\sum X_i^2 \sim \chi^2(n)$ with PDF involving $\Gamma(n/2)$.
\end{remark}

\begin{remark}[Riemann Zeta Function]
The functional equation relates $\zeta(s)$ to $\zeta(1-s)$ via the gamma function:
\[
\zeta(s) = 2^s \pi^{s-1} \sin\left(\frac{\pi s}{2}\right) \Gamma(1-s) \zeta(1-s).
\]
\end{remark}

\section{Frequently Asked Questions}

\begin{itemize}
\item \textbf{Q: Why is $\Gamma(n) = (n-1)!$ instead of $\Gamma(n) = n!$?} \\
The convention $\Gamma(n) = (n-1)!$ comes from the integral definition. We have $\Gamma(1) = 1 = 0!$, $\Gamma(2) = 1 = 1!$, $\Gamma(3) = 2 = 2!$, so $\Gamma(n) = (n-1)!$. This makes the recurrence $\Gamma(z+1) = z\Gamma(z)$ work cleanly.

\item \textbf{Q: What is $\Gamma(1/2)$?} \\
$\Gamma(1/2) = \sqrt{\pi} \approx 1.77245$. This follows from the Gaussian integral $\int_{-\infty}^{\infty} e^{-x^2} dx = \sqrt{\pi}$ via the substitution $t = x^2$.

\item \textbf{Q: What are the poles of the Gamma function?} \\
$\Gamma(z)$ has simple poles at $z = 0, -1, -2, -3, \ldots$ with residue $\frac{(-1)^n}{n!}$ at $z = -n$.

\item \textbf{Q: How is the Gamma function used in probability?} \\
The Gamma distribution has PDF $f(x;\alpha, \beta) = \frac{\beta^\alpha}{\Gamma(\alpha)} x^{\alpha-1} e^{-\beta x}$. The Beta distribution uses the ratio $\frac{\Gamma(\alpha)\Gamma(\beta)}{\Gamma(\alpha+\beta)}$.

\item \textbf{Q: What is the relationship between $\Gamma$ and the factorial?} \\
The Gamma function extends the factorial to complex numbers (except non-positive integers). $\Gamma(n+1) = n!$ for non-negative integers $n$.

\item \textbf{Q: How is the Beta function defined?} \\
The Beta function is $B(x,y) = \frac{\Gamma(x)\Gamma(y)}{\Gamma(x+y)} = \int_0^1 t^{x-1}(1-t)^{y-1} dt$. It appears in Bayesian statistics and combinatorics.

\item \textbf{Q: What is the connection to the Riemann zeta function?} \\
The functional equation $\zeta(s) = 2^s \pi^{s-1} \sin(\pi s/2) \Gamma(1-s) \zeta(1-s)$ connects values of $\zeta$ at $s$ and $1-s$ through the gamma function.

\item \textbf{Q: Can $\Gamma$ be computed for negative numbers?} \\
Yes, using the reflection formula $\Gamma(z)\Gamma(1-z) = \pi/\sin(\pi z)$. For negative non-integer $z$, this gives a finite value.
\end{itemize}

\section{Exercises}

\begin{exercise}
Prove that $\Gamma(1/2) = \sqrt{\pi}$ using the Gaussian integral.
\end{exercise}

\begin{exercise}
Use the reflection formula to show $\Gamma(1/4)\Gamma(3/4) = \pi\sqrt{2}$.
\end{exercise}

\begin{exercise}
Derive Stirling's formula from the integral definition using Laplace's method.
\end{exercise}

\begin{exercise}
Show that $\psi(n) = -\gamma + \sum_{k=1}^{n-1} \frac{1}{k}$ for positive integers $n$.
\end{exercise}

\begin{exercise}[Computational]
Implement the Lanczos approximation for computing $\Gamma(z)$ and compare accuracy with the built-in `math.gamma`.
\end{exercise}

\begin{exercise}
Verify the duplication formula $\Gamma(z)\Gamma(z+1/2) = 2^{1-2z}\sqrt{\pi}\Gamma(2z)$ for $z = 1/2, 1, 3/2$.
\end{exercise}

\begin{exercise}
Prove the Beta function identity $B(x,y) = \frac{\Gamma(x)\Gamma(y)}{\Gamma(x+y)}$ using the substitution $t = \frac{u}{u+v}$.
\end{exercise}

\begin{exercise}
Use the Gauss multiplication formula to show $\Gamma(1/3)\Gamma(2/3) = \frac{2\pi}{\sqrt{3}}$.
\end{exercise}

\begin{exercise}
Compute $\Gamma(1/4)$ numerically and verify that $\Gamma(1/4) \approx 3.6256$.
\end{exercise}


\section{Special Values and Identities}

\begin{theorem}[Gamma at Half-Integers]
For any positive integer $n$:
\[
\Gamma\left(n + \frac{1}{2}\right) = \frac{(2n)!}{4^n n!} \sqrt{\pi} = \frac{(2n-1)!!}{2^n} \sqrt{\pi}.
\]
\end{theorem}

\begin{proof}
Using the duplication formula with $z = 1/2$:
\[
\Gamma\left(\frac{1}{2}\right)\Gamma(1) = 2^{0}\sqrt{\pi}\,\Gamma(1) \implies \Gamma\left(\frac{1}{2}\right) = \sqrt{\pi}.
\]
Then applying the recurrence relation repeatedly gives the general formula.
\end{proof}

\begin{corollary}
\[
\Gamma\left(\frac{1}{2}\right) = \sqrt{\pi}, \quad \Gamma\left(\frac{3}{2}\right) = \frac{\sqrt{\pi}}{2}, \quad \Gamma\left(\frac{5}{2}\right) = \frac{3\sqrt{\pi}}{4}, \quad \Gamma\left(\frac{7}{2}\right) = \frac{15\sqrt{\pi}}{8}.
\]
\end{corollary}

\begin{theorem}[Gamma at Rational Points]
\[
\Gamma\left(\frac{1}{3}\right)\Gamma\left(\frac{2}{3}\right) = \frac{2\pi}{\sqrt{3}}, \qquad \Gamma\left(\frac{1}{4}\right)\Gamma\left(\frac{3}{4}\right) = \pi\sqrt{2}.
\]
\end{theorem}

\section{Gauss's Formula}

\begin{theorem}[Gauss's Digamma Theorem]
For any positive integers $p$ and $q$ with $p < q$:
\[
\psi\left(\frac{p}{q}\right) = -\gamma - \log(2q) - \frac{\pi}{2}\cot\left(\frac{p\pi}{q}\right) + 2\sum_{n=1}^{\lfloor(q-1)/2\rfloor} \cos\left(\frac{2\pi np}{q}\right)\log\left(\sin\left(\frac{n\pi}{q}\right)\right).
\]
\end{theorem}

\section{Euler's Constant Formulas}

\begin{theorem}[Integral for $\gamma$]
\[
\gamma = -\int_0^\infty e^{-t}\log t\, dt.
\]
\end{theorem}

\begin{proof}
From the definition $\Gamma(z) = \int_0^\infty t^{z-1} e^{-t} dt$, differentiating at $z = 1$:
\[
\Gamma'(1) = \int_0^\infty \log t \cdot e^{-t} dt = -\gamma.
\]
\end{proof}

\begin{theorem}[Euler's Series]
\[
\gamma = \sum_{k=1}^\infty \left(\frac{1}{k} - \log\left(1 + \frac{1}{k}\right)\right).
\]
\end{theorem}

\begin{theorem}[Binet's Second Formula]
\[
\log\Gamma(z) = \left(z - \frac{1}{2}\right)\log z - z + \frac{1}{2}\log(2\pi) + \int_0^\infty \left(\frac{1}{e^t - 1} - \frac{1}{t} + \frac{1}{2}\right)\frac{e^{-zt}}{t} dt.
\]
\end{theorem}

\section{Numerical Properties}

The Gamma function satisfies several useful bounds:

\begin{theorem}[ Gautschi's Inequality]
For $x > 0$ and $0 < s < 1$:
\[
x^{1-s} < \frac{\Gamma(x+1)}{\Gamma(x+s)} < \left(x + \frac{1}{2}\right)^{1-s}.
\]
\end{theorem}

\begin{theorem}[Kershaw's Inequality]
For $x > 0$ and $0 < s < 1$:
\[
\left(x + \frac{s}{2}\right)^{1-s} < \frac{\Gamma(x+1)}{\Gamma(x+s)} < \left(x + \sqrt{x^2 + s}\right)^{1-s}.
\]
\end{theorem}

\section{Higher Derivatives}

\begin{theorem}[Polygamma Functions]
The $n$-th derivative of $\Gamma(z)$ can be expressed using polygamma functions:
\[
\psi^{(n)}(z) = \frac{d^{n+1}}{dz^{n+1}} \log\Gamma(z) = (-1)^{n+1} n! \sum_{k=0}^\infty \frac{1}{(z+k)^{n+1}}.
\]
\end{theorem}

\begin{corollary}
\[
\psi^{(n)}(1) = (-1)^{n+1} n! \zeta(n+1).
\]
\end{corollary}

\begin{itemize}
    \item $\psi^{(0)}(1) = -\gamma$
    \item $\psi^{(1)}(1) = \zeta(2) = \frac{\pi^2}{6}$
    \item $\psi^{(2)}(1) = -2\zeta(3)$
    \item $\psi^{(3)}(1) = 6\zeta(4) = \frac{\pi^4}{15}$
\end{itemize}

\section{Connections to Other Functions}

\subsection{Beta Function}

The Beta function $B(x,y)$ is intimately related to the Gamma function:
\[
B(x,y) = \frac{\Gamma(x)\Gamma(y)}{\Gamma(x+y)}.
\]

This identity is used extensively in probability theory and statistics.

\subsection{Wallis Product}

Setting $z = 1/2$ in the reflection formula and using $\Gamma(1/2) = \sqrt{\pi}$ yields:
\[
\frac{\pi}{2} = \prod_{n=1}^\infty \frac{4n^2}{4n^2 - 1} = \frac{2}{1} \cdot \frac{2}{3} \cdot \frac{4}{3} \cdot \frac{4}{5} \cdot \frac{6}{5} \cdot \frac{6}{7} \cdots
\]

\subsection{Weierstrass Product}

The Weierstrass product for $\Gamma(z)$ connects it to the Riemann zeta function:
\[
\frac{1}{\Gamma(z)} = z e^{\gamma z} \prod_{n=1}^\infty \left(1 + \frac{z}{n}\right) e^{-z/n}.
\]

Taking logarithms and differentiating yields the digamma function:
\[
\psi(z) = -\gamma - \frac{1}{z} + \sum_{n=1}^\infty \left(\frac{1}{n} - \frac{1}{n+z}\right).
\]
